\documentclass[a4paper]{article}

\usepackage[fontset=fandol]{ctex}
\usepackage{amsmath, amssymb}
\usepackage{graphicx}
\usepackage[margin=2.5cm]{geometry}
\usepackage[most]{tcolorbox}
\usepackage{etoolbox}
\usepackage{listings}
\usepackage{hyperref}
\usepackage{booktabs}
\usepackage{subcaption}
\usepackage{float}
\usepackage{tikz}
\IfFileExists{pgfplots.sty}{
    \usepackage{pgfplots}
    \pgfplotsset{compat=1.18}
}{}

\newtcolorbox{knowledgebox}[1]{
    enhanced, colback=blue!5!white, colframe=blue!75!black, colbacktitle=blue!75!black,
    coltitle=white, fonttitle=\bfseries, title=#1, attach boxed title to top left={yshift=-2mm, xshift=2mm},
    boxrule=1pt, sharp corners
}

\newtcolorbox{importantbox}[1]{
    enhanced, colback=yellow!10!white, colframe=yellow!80!black, colbacktitle=yellow!80!black,
    coltitle=black, fonttitle=\bfseries, title=#1, sharp corners
}

\newtcolorbox{warningbox}[1]{
    enhanced, colback=red!5!white, colframe=red!75!black, colbacktitle=red!75!black,
    coltitle=white, fonttitle=\bfseries, title=#1, sharp corners
}

\lstset{
    language=Python,
    basicstyle=\ttfamily\small,
    keywordstyle=\color{blue},
    stringstyle=\color{red!60!black},
    commentstyle=\color{green!60!black},
    breaklines=true,
    frame=single,
    numbers=left,
    numberstyle=\tiny\color{gray},
    captionpos=b,
    extendedchars=false
}

\newcommand{\notetitle}{CS224R Lecture 12: Multi-Task and Goal-Conditioned RL}
\newcommand{\noteauthors}{Codex}
\newcommand{\notedate}{2026-04-13}
\newcommand{\videochannel}{Stanford CS224R}
\newcommand{\videopublishdate}{2025-05-09}
\newcommand{\videoduration}{约 01:10:29}
\newcommand{\videourl}{https://www.youtube.com/watch?v=qNdsI_4AQJw}
\newcommand{\videocoverpath}{cover.jpg}

\begin{document}

\begin{titlepage}
\centering
{\Large 课程笔记\par}
\vspace{1.2cm}
{\huge\bfseries \notetitle\par}
\vspace{0.8cm}
{\large \noteauthors\par}
\vspace{0.3cm}
{\large \notedate\par}
\vspace{1.2cm}

\includegraphics[width=0.82\textwidth,height=0.45\textheight,keepaspectratio]{\videocoverpath}\par

\vfill
\begin{tcolorbox}[width=0.9\textwidth, colback=black!2!white, colframe=black!60, sharp corners]
\textbf{视频作者/频道}：\videochannel\par
\textbf{发布日期}：\videopublishdate\par
\textbf{视频时长}：\videoduration\par
\textbf{视频链接}：\href{\videourl}{\nolinkurl{\videourl}}
\end{tcolorbox}
\end{titlepage}

\tableofcontents
\newpage

\section{先补上上一讲未讲完的一半：如何用模型做 policy optimization}

这讲虽然题目是 multi-task 与 goal-conditioned RL，但开头并不是直接跳入多任务，而是先承接上一讲，补完 model-based policy optimization 的另一种用法。原因很合理：如果 planning 只适合短 horizon、测试时又太贵，那么 learned model 更现实的作用往往不是直接在线规划，而是帮助训练策略。

\subsection{为什么 imagined rollouts 要短、而且要从真实状态出发}

讲者首先回顾了 planning with learned models 的优缺点：优点是简单，可以直接换目标和奖励；缺点是测试时算力重，而且对长 horizon 不现实，因为既贵又会因模型误差累积而失真。于是自然的问题是：\textbf{能不能把模型的价值转移到训练阶段，而不是测试阶段}？

答案是肯定的。model-based policy optimization 的关键直觉，是用 learned model 生成合成数据，补充真实环境里拿到的轨迹。但讲者反复强调，不能从初始状态一口气生成很长 imagined trajectory；更好的做法，是从 replay buffer 里的很多真实状态出发，只生成\textbf{短的} synthetic rollouts。

\begin{figure}[H]
\centering
\includegraphics[width=0.9\textwidth]{pdf_pages/page-07.png}
\caption{model-based policy optimization 的完整流程：环境数据训练模型，模型从真实状态附近生成短 imagined rollouts，再反过来更新策略。}
\end{figure}

这种做法的好处有两个。第一，它让合成轨迹始终锚定在真实数据邻域，不容易完全漂走。第二，它把模型误差控制在短 horizon 内，更适合作为局部覆盖增强，而不是全盘替代环境。

\subsection{什么时候应该考虑 model-based RL}

课程没有把 model-based RL 当作默认更高级的方法，而是给出一套非常明确的取舍标准：如果模型容易学，它能极大提升样本效率，并且还能在没有 reward 标签时用自监督方式训练；如果模型难学，它就会变成新的误差源、新的超参数源和新的计算负担。

\begin{figure}[H]
\centering
\includegraphics[width=0.88\textwidth]{pdf_pages/page-08.png}
\caption{是否使用 model-based RL，核心取决于模型是否足够好学，而不是它在概念上是否 ``更聪明''。}
\end{figure}

这一点和后半节的多任务主题其实是相通的：无论是 learned dynamics 还是 shared multi-task structure，真正有用的前提都是存在可复用、可泛化的结构。如果底层结构本身不存在，共享和迁移就无从谈起。

\subsection{本章小结}

Lecture 12 的开头先把上一讲的故事补完整了。规划不是模型的唯一用途；在很多现实系统中，模型更像一个训练期的数据放大器。只要 imagined rollouts 足够短、足够贴近真实分布，它们就能为策略学习提供额外的局部探索信号。

\section{为什么我们想学一个 generalist policy}

\subsection{multi-task RL 的动机不只是 ``一模多用''}

讲者随后把话题转到 multi-task RL。表面上看，多任务的目标很直观：不想只学一个会做单一事情的策略，而是想学一个 generalist policy，能处理多个目标、多个任务、多个用户或多个环境。例子横跨语言助手、腿足机器人、移动机械臂、推荐系统和游戏代理。

\begin{figure}[H]
\centering

\includegraphics[width=0.9\textwidth]{pdf_pages/page-11.png}
\caption{multi-task RL 关注的是 generalist policy：一个策略需要在任务标识变化时展现出不同能力。}
\end{figure}

但课程真正强调的不是 ``方便''，而是\textbf{数据效率}。如果多个任务共享结构，例如语言助手共享语法与网页操作模式，腿足机器人共享平衡能力，不同用户推荐共享偏好簇结构，那么多个任务联合训练，就可能把单任务本来需要反复独立学习的知识 amortize 掉。

\subsection{generalist 往往比 specialist 更稳}

讲者给出的第二个动机更值得注意：经验上，generalist systems 往往不只是覆盖面更广，还常常比 specialist 更可靠、更高性能。这并非逻辑必然，但在近年 AI 系统里屡次出现。课程用它来说明：多任务并不只是工程上的合并，同样可能是性能提升的来源。

\begin{knowledgebox}{这节课的关键转折}
从这里开始，课程把 ``如何学一个任务'' 变成 ``如何在共享结构中同时学多个任务''。这意味着状态、动作、奖励和任务标识之间的关系必须重新形式化，否则多任务只会停留在口号层面。
\end{knowledgebox}

\subsection{本章小结}

multi-task RL 的根本目标不是省模型数量，而是利用任务之间的共享结构提升样本效率和泛化。只有把这种共享写进问题形式里，后续的 task conditioning、relabeling 和 goal-conditioned RL 才有意义。

\section{如何把 ``任务'' 写进 MDP}

\subsection{把不同任务当成不同 MDP}

课程给出的第一种形式化方式，是把每个任务 $i$ 看成一个独立的 MDP：
\[
\mathcal{M}_i=(\mathcal{S}_i,\mathcal{A}_i,p_i(s_1),p_i(s'|s,a),r_i(s,a)).
\]
这里不同任务可以共享某些组成部分，也可以在状态空间、动作空间、动力学、初始状态分布、奖励函数上彼此不同。讲者特别提醒，这种定义比日常语言里的 ``任务'' 更一般，因为它允许我们把任何 MDP 差异都视为任务差异。

\begin{figure}[H]
\centering
\includegraphics[width=0.86\textwidth]{pdf_pages/page-13.png}
\caption{把任务形式化为一组不同的 MDP：状态、动作、动力学、初始状态分布与奖励函数都可能因任务而变化。}
\end{figure}

这个定义的价值，在于它直接告诉我们多任务共享可能发生在哪一层。比如不同 locomotion skill 往往共享动力学但奖励不同；不同 garment manipulation 任务则可能连动力学和初始状态分布都变；不同用户推荐的任务，奖励偏好会变，但系统交互接口可能保持相同。

\subsection{把 task identifier 放进状态里}

如果多个任务至少共享某个统一的表示空间，课程就给出第二种更便于算法实现的写法：把任务标识 $z_i$ 直接并入状态，得到聚合后的 MDP。此时状态写作
\[
s=(\bar{s}, z_i),
\]
相应的奖励与动力学也写成依赖任务标识的形式：
\[
r((\bar{s}, z_i), a)=r_i(\bar{s}, a), \qquad
p((\bar{s}', z_i)\mid (\bar{s}, z_i), a)=p_i(\bar{s}'\mid \bar{s}, a).
\]

对于 goal-conditioned RL，则常把任务标识具体化为目标状态 $g$。这时一个常见奖励形式是
\[
r(\bar{s}, g)=-d(\bar{s}, g),
\]
也就是用当前状态与目标之间的距离来定义回报。距离既可以是 Euclidean $\ell_2$，也可以是稀疏 $0/1$ 成功信号。

\begin{figure}[H]
\centering
\includegraphics[width=0.84\textwidth]{pdf_pages/page-17.png}
\caption{把 task identifier 并入状态后，多任务问题就可以转写为一个扩展状态空间上的标准 RL 问题。}
\end{figure}

讲者明确指出：一旦任务标识成为状态的一部分，\textbf{标准 RL 算法其实就可以直接用}。这为后面所有 multi-task imitation、multi-task RL、goal-conditioned RL 提供了统一入口。

\subsection{本章小结}

这一部分完成了多任务问题的数学重写。课程先把不同任务视为不同 MDP，再说明只要能把任务标识嵌入状态，就能把多任务问题重新变成一个 ``标准'' 的 RL 问题。后续的关键便转向：这个 task identifier 在实践中怎么表示、怎么使用，以及数据能否在任务之间共享。

\section{task conditioning 在实践中怎么做}

\subsection{从多任务模仿学习开始}

课程先从 multi-task imitation learning 讲起，因为它最容易看清 task conditioning 的本质。对于单任务 BC，策略只看状态；到了多任务版本，策略要看状态和任务标识。一个简单但很重要的训练细节，是在每个 minibatch 里做分任务的 stratified sampling，降低梯度方差，避免数据量大的任务淹没数据量小的任务。

\subsection{任务标识可以是 one-hot，也可以是语言或示范}

讲者没有把 task identifier 限制成 one-hot ID。相反，课程明确展示了多种可能的 task conditioning 形式：one-hot 任务编号、语言描述、示范视频、目标状态，甚至在更现代的系统里，任务标识可以作为 prompt 输入给大模型策略。

\begin{figure}[H]
\centering
\includegraphics[width=0.94\textwidth]{pdf_pages/page-19.png}
\caption{task conditioning 的实践形态并不单一：视觉输入、示范视频、语言描述都可以共同构成任务嵌入。}
\end{figure}

这部分非常贴近当代 VLA 与 language-conditioned robotics 的做法。任务不再是一个抽象编号，而是一个可学习的上下文变量。讲者还提到，如果任务本质上是 ``不同人的偏好''，那么更好的表示甚至可能是几条该用户喜欢或不喜欢的例子，这就有点接近 in-context learning 的味道了。

\subsection{多任务算法只是在单任务算法外再接一个条件}

从算法视角看，multi-task RL 并没有完全改写 RL 的基本对象，而是在 policy 和 Q-function 上显式加入任务条件：
\[
\pi_\theta(a\mid \bar{s}) \longrightarrow \pi_\theta(a\mid \bar{s}, z_i),
\]
\[
Q_\phi(\bar{s}, a) \longrightarrow Q_\phi(\bar{s}, a, z_i).
\]
这说明一个很重要的事实：\textbf{多任务 RL 的难点不在公式是否复杂，而在任务之间哪些数据真能共享，哪些共享会制造伪相关。}

\subsection{本章小结}

task conditioning 的本质，是把 ``当前应该做哪件事'' 变成策略输入的一部分。形式上这看似只是多加一个条件变量，实际上它决定了模型能否正确对齐不同任务语义，也决定了后续跨任务复用数据是否可能。

\section{跨任务共享数据：从 hindsight relabeling 到 goal-conditioned RL}

\subsection{为什么可以把一条轨迹重解释成别的任务}

课程接着讨论一个极有启发性的想法：如果一条轨迹本来是为了任务 $z_1$ 采集的，但它在执行过程中顺便完成了任务 $z_2$，那这条数据能不能也拿来训练任务 $z_2$？这就是 hindsight relabeling 的基本直觉。

\begin{figure}[H]
\centering
\includegraphics[width=0.9\textwidth]{pdf_pages/page-24.png}
\caption{hindsight relabeling 的直觉：一条轨迹若意外完成了另一个任务，也可以被重新标注后再次利用。}
\end{figure}

在一般 multi-task RL 中，重标注可以写成：先收集
\[
\mathcal{D}_k=\{(s_{1:T},a_{1:T},z_i,r_{1:T})\},
\]
然后把它改写为另一个任务 $z_j$ 下的版本
\[
\mathcal{D}'_k=\{(s_{1:T},a_{1:T},z_j,r'_{1:T})\},
\qquad
r'_t=r_j(s_t).
\]
只要新的任务奖励可以在事后计算，这种 relabeling 就让同一条经验服务于多个任务。

\subsection{goal-conditioned RL 是最自然的 relabeling 场景}

goal-conditioned RL 是 hindsight relabeling 最经典也最自然的应用。因为在目标到达任务里，轨迹里出现过的某个后续状态，本身就可以被重新解释为 ``原来这条轨迹的目标就是到这里''。

\begin{figure}[H]
\centering
\includegraphics[width=0.9\textwidth]{pdf_pages/page-26.png}
\caption{goal-conditioned RL with hindsight relabeling：把轨迹终点或中间状态改写成目标，就能把失败探索变成成功样本。}
\end{figure}

课上给出的算法流程非常清楚：

\begin{enumerate}
\item 收集一条以目标 $g$ 为条件的轨迹；
\item 存进 replay buffer；
\item 事后把终点状态 $s_T$ 或轨迹中的其他状态改写成新的目标；
\item 用新目标重新计算奖励，再把重标注后的轨迹也放回 buffer；
\item 用 off-policy RL 从混合 buffer 中更新策略。
\end{enumerate}

它最核心的收益，在于\textbf{把原本稀疏、经常失败的探索，转化为事后看来成功的训练样本}。这就是 HER, Hindsight Experience Replay 的经典思想。

\subsection{但不是所有任务都能随便 relabel}

讲者也明确指出，relabeling 并非普适魔法。它成立需要若干前提：

\begin{itemize}
\item 不同任务至少要共享兼容的状态空间与动作空间；
\item 奖励函数形式必须已知且可事后评估；
\item 动力学必须与新的任务解释相容；
\item 通常需要 off-policy 算法，因为 replay buffer 里会混入许多并非按当前任务采样的经验。
\end{itemize}

\begin{importantbox}{什么时候 relabeling 最有价值}
当任务本身稀疏、探索难，而且 ``失败轨迹'' 中蕴含大量潜在成功子目标时，relabeling 的价值最大。goal-conditioned robotic manipulation 正是这种结构最明显的场景。
\end{importantbox}

\subsection{本章小结}

跨任务共享数据并不是简单地把所有轨迹混在一起，而是要利用任务结构进行有约束的重解释。hindsight relabeling 的妙处，在于它利用了任务之间的共同动力学，把 ``没完成原目标'' 的轨迹，转化成 ``完成了另一个可行目标'' 的样本，从而显著缓解探索困难。

\section{总结与延伸}

Lecture 12 实际上做了两件彼此呼应的事。前半段说明 learned model 可以作为 imagined rollout generator，帮助训练策略；后半段说明 multi-task structure 可以作为共享与重标注的依据，帮助同一批经验服务于多个任务。两者背后的思想是一致的：\textbf{如果环境里存在可复用的结构，就应该想办法把结构转换成数据效率}。

在多任务部分，课程完成了一条非常完整的逻辑链。首先，多任务的真正动机是样本效率与泛化，而不是模型数量管理。其次，任务可以被写成不同 MDP，也可以通过把 task identifier 并入状态，转化为扩展状态空间上的单一 MDP。再次，task conditioning 在实践中可以是 one-hot、语言、示范视频或目标状态。最后，若任务之间共享动力学且奖励可重算，就能使用 hindsight relabeling，把本来只对应单一任务的经验复用到其他任务上。

从研究角度看，这节课留下的最值得继续追问的问题有三个。第一，task representation 应该如何学习，才能真正捕捉任务间共享结构，而不是只记住任务 ID。第二，relabeling 什么时候会引入错误归因，尤其在任务边界不清或奖励函数近似时。第三，随着多任务系统越来越接近 language-conditioned agents，task conditioning 是否会进一步演化为更一般的上下文学习或 preference conditioning。对 CS224R 来说，这正是从经典 RL 走向 modern generalist agent 的关键过渡。

\end{document}
